% Slajd 7 – dostępność cyfrowa: wykaz funkcji dostępnych i elementów do dalszych prac (pkt 5 i 6 wyzwania).
\begin{frame}{Dostępność cyfrowa: cel WCAG 2.2 AA}
  \begin{columns}[T,onlytextwidth]
    \begin{column}{0.62\textwidth}
      \begin{block}{Jest w~prototypie (sprawdzone w~głównym scenariuszu)}
        \footnotesize
        \begin{itemize}
          \item \textbf{Klawiatura:} każda akcja to przycisk lub link; link „Przejdź do treści”; widoczny fokus,
                który nie ginie przy odświeżaniu list co 15~s; Esc zamyka panele.
          \item \textbf{Czytnik ekranu:} landmarki, nagłówki, etykiety ARIA, krótkie komunikaty „live”;
                opis mapy odsyła do listy; tytuł okna zmienia się z~widokiem.
          \item \textbf{Mapa jako tekst:} listy Zgłoszenia, Miejsca i~najbliższych obiektów warstw;
                waga problemu kształtem i~podpisem, nie tylko kolorem.
          \item \textbf{Kontrast i~czytelność:} tokeny 4,5:1 (tekst) i~3:1 (kontrolki) sprawdzane testem;
                tryb nocny; czcionka Atkinson Hyperlegible; 4~języki.
          \item \textbf{Formularze i~dotyk:} błędy w~\texttt{role="alert"} z~nazwą pola, „Dalej” nie jest blokowane;
                cele $\geq$~44~px; układ działa do 400\,\% powiększenia.
        \end{itemize}
      \end{block}
    \end{column}
    \begin{column}{0.36\textwidth}
      \begin{alertblock}{Do dalszych prac}
        \footnotesize
        \begin{itemize}
          \item Audyt z~NVDA / VoiceOver i~testy z~użytkownikami na wózkach i~słabowidzącymi.
          \item Klawiatura po znacznikach mapy (Leaflet); alternatywa dla ustawiania pinezki
                przesuwaniem mapy (2.5.7).
          \item Panele właściciela i~administratora w~tym samym standardzie.
          \item Deklaracja dostępności i~raport zgodności.
        \end{itemize}
      \end{alertblock}
      \vspace{0.5mm}
      \muted{Iteracje „WCAG 2.2 A” i~„WCAG 2.2 AA” są objęte testami \texttt{AccessibilityTests}.}
    \end{column}
  \end{columns}

  \note[item]{Pokazać na żywo: Tab przez listę zgłoszeń, Enter otwiera szczegóły, Esc wraca; przełączyć tryb nocny.}
  \note[item]{„Tekstowa alternatywa mapy” to wymaganie wprost z~wyzwania – u~nas lista obok mapy jest pierwszorzędnym widokiem, nie dodatkiem.}
  \note[item]{Kontrast: test liczy współczynniki z~tokenów w~\texttt{style.css} dla obu motywów; separatory mają własny token, obwódki kontrolek i~symbole warstw mają 3:1.}
  \note[item]{Nie mówić „jesteśmy zgodni z~WCAG”; mówić „główny scenariusz spełnia kryteria A/AA, które sprawdziliśmy, a~tu jest lista braków”.}
\end{frame}
